\documentclass[11pt,letterpaper]{article}
\usepackage{ad42010-suite}
\begin{document}
\DocFront{03}{Architecture Viewpoint Catalog}{Viewpoint specifications, model kinds, legends and methods}{A reusable catalog of eight original software-intensive-system viewpoint candidates with explicit metadata required for selecting and governing architecture views.}
\InterpretationKey
\section{Catalog status and specification contract}
\begin{Requirement}\FlagReq{} A viewpoint specification must cover associated stakeholder perspectives, framed concerns, related aspects, typical stakeholders, model kinds and legends used for views, correspondence methods and information sources. Interpretive legends need guidance for the view components they document. \Source{8.1}\end{Requirement}
\begin{Guidance}The following entries are \emph{proposed viewpoint specifications}, not mandated viewpoints. Their model kinds and methods are examples. Before use, assign a controlled version, ensure referenced model kinds meet 8.2, adapt to actual stakeholders and concerns and register the selected viewpoint in the AD. Annex B provides informative viewpoint-documentation guidance.\end{Guidance}
\section{Reusable viewpoint specification form}
\begin{longtable}{p{0.30\linewidth}p{0.60\linewidth}}\toprule\textbf{Field}&\textbf{Specification content}\\\midrule\endhead
Name / identifier / version&Unique label, edition and ownership, with applicability.\\
Overview / excluded questions&Viewpoint objective and the concerns it is not intended to frame.\\
Stakeholder perspectives&Named perspectives associated with viewpoint; link to local perspective register.\\
Typical stakeholders&Roles and classes likely to hold the framed concerns.\\
Concerns&One or more questions or evaluation targets explicitly framed.\\
Aspects&Characteristics of EoI relevant to these concerns.\\
Model kinds and legends&Referenced model-kind specs, legend specifications and governing conventions.\\
Correspondence methods&Rules for cross-component or cross-view relationships and consistency.\\
View methods&Construction, interpretation and analysis techniques; recommended record.\\
Sources&Publications, standards, domain practices or internally approved references.\\
Example views&Illustrative view names, view-component shapes and issues to capture.\\\bottomrule\end{longtable}
\Source{8.1; Annex B.2 (informative)}
\section{Candidate architecture viewpoints}
\subsection{VP-CXT: Entity Context viewpoint (version 1.0)}
\textbf{Purpose}: frame the EoI's boundary, interactions and environmental constraints. \textbf{Perspectives}: business, operational and acquisition. \textbf{Typical stakeholders}: product owners, external integrators, operators and risk reviewers.\par
\textbf{Concerns}: Where is the EoI boundary? What outside entities and rules influence it? Which interfaces carry key interactions? \textbf{Aspects}: structural, functional and external context.\par
\textbf{Model kinds}: \texttt{MK-CONTEXT/1}; entity-interaction graph with typed vertices/edges. \textbf{Legends}: interface categories, internal/external visual markers, trust-boundary label and arrow meaning. \textbf{Correspondence method}: \texttt{CM-EXT/1}, each external dependency in the context component must have a linked interface or documented exception.\par
\textbf{View methods}: stakeholder interview, interface inventory, edge review and boundary interrogation. \textbf{Sources}: project scope baseline, API inventory and stakeholder records; conventions are locally specified here. \textbf{Example views}: partner integration context and jurisdictional boundary.\par
\subsection{VP-FUN: Functional Behaviour viewpoint (version 1.0)}
\textbf{Purpose}: explain responsibilities, event flows, state transitions and failure behaviour. \textbf{Perspectives}: operational and development. \textbf{Stakeholders}: engineers, operators, testers.\par
\textbf{Concerns}: What happens along critical scenarios? What work is synchronous? How do failure modes recover? \textbf{Aspects}: functional and behavioural.\par
\textbf{Model kinds}: \texttt{MK-SEQUENCE/1} (participant interactions), \texttt{MK-STATE/1} (state transitions). \textbf{Legends}: activation, message direction, guards, timeout notation. \textbf{Correspondence method}: \texttt{CM-EVENT/1}, interface calls in runtime behaviour resolve to interfaces declared in context/deployment records.\par
\textbf{View methods}: scenario walkthrough, negative-path enumeration and trace-to-test analysis. \textbf{Sources}: documented sequence/state modelling conventions, functional requirements. \textbf{Example view}: transaction failure and retry.\par
\subsection{VP-INFO: Information viewpoint (version 1.0)}
\textbf{Purpose}: describe information semantics, ownership and movement. \textbf{Perspectives}: data, business and assurance. \textbf{Stakeholders}: data owners, engineers, privacy and security reviewers.\par
\textbf{Concerns}: Which data are critical or sensitive? Where do data originate, traverse and persist? Who is responsible for lifecycle handling? \textbf{Aspects}: informational and structural.\par
\textbf{Model kinds}: \texttt{MK-DATAFLOW/1}, \texttt{MK-DATADICTIONARY/1}. \textbf{Legends}: data classification codes, store/flow symbols and lineage directions. \textbf{Correspondence method}: \texttt{CM-DATA/1}, each sensitive flow links to a classified data object and a location or interface.\par
\textbf{View methods}: source-to-sink tracing, classification verification, retention analysis. \textbf{Sources}: approved data inventory and interface contracts. \textbf{Example view}: personal-data transfer chain.\par
\subsection{VP-DEP: Deployment and Allocation viewpoint (version 1.0)}
\textbf{Purpose}: make runtime placement and resource dependencies understandable. \textbf{Perspectives}: physical/technical, operational. \textbf{Stakeholders}: platform engineer, reliability engineer, infrastructure owner.\par
\textbf{Concerns}: What runs where? Where do fault and trust boundaries lie? How do capacity constraints affect placement? \textbf{Aspects}: structural, physical and programmatic.\par
\textbf{Model kinds}: \texttt{MK-DEPLOY/1}, a typed allocation graph. \textbf{Legends}: compute node, zone, cluster, network path and failover marker. \textbf{Correspondence method}: \texttt{CM-ALLOC/1}, every deployable runtime element maps to at least one declared deployment target.\par
\textbf{View methods}: environment inventory reconciliation, availability-zone walkthrough, allocation checks. \textbf{Sources}: infrastructure manifests, environment topology and service catalog. \textbf{Example view}: production cluster placement.\par
\subsection{VP-SEC: Security and Trust viewpoint (version 1.0)}
\textbf{Purpose}: address abuse cases, trust boundaries and security control placement. \textbf{Perspectives}: assurance, operations and development. \textbf{Stakeholders}: application security engineer, security architect, risk owner, operator.\par
\textbf{Concerns}: Which interactions cross trust boundaries? Which components are privileged? How are threats mitigated or accepted? \textbf{Aspects}: informational, structural and behavioural.\par
\textbf{Model kinds}: \texttt{MK-THREAT/1} (threat/asset/control relationships) and \texttt{MK-DATAFLOW/1}. \textbf{Legends}: threat IDs, security boundaries, control status and evidence source. \textbf{Correspondence method}: \texttt{CM-CTRL/1}, each material identified threat has a treatment decision and owner.\par
\textbf{View methods}: threat elicitation, attack-path tracing and design control verification. \textbf{Sources}: authorized threat model, system risk register and security requirements; threat methods such as STRIDE are optional external practices. \textbf{Example view}: authentication and authorization trust boundary.\par
\subsection{VP-OPS: Operations and Resilience viewpoint (version 1.0)}
\textbf{Purpose}: expose operational ownership, recovery and observability. \textbf{Perspectives}: operations, support, service management. \textbf{Stakeholders}: on-call engineer, operator, incident commander, service owner.\par
\textbf{Concerns}: What signals reveal degradation? How are failures contained and recovery performed? Who acts on an incident? \textbf{Aspects}: behavioural, functional and programmatic.\par
\textbf{Model kinds}: \texttt{MK-FAILURE/1}, a dependency/failure propagation graph; \texttt{MK-STATE/1}. \textbf{Legend}: severity, observability coverage, recovery path and runbook ownership. \textbf{Correspondence method}: \texttt{CM-RUN/1}, each critical service node traces to an operating procedure and monitoring signal.\par
\textbf{View methods}: tabletop exercise, fault injection review, SLO budget analysis. \textbf{Sources}: service-level objectives, incident records and production runbooks. \textbf{Example view}: provider outage containment.\par
\subsection{VP-EVO: Evolution and Change viewpoint (version 1.0)}
\textbf{Purpose}: document transitions, dependencies and migration trade-offs. \textbf{Perspectives}: roadmap, acquisition and development. \textbf{Stakeholders}: architect, engineering lead, portfolio owner, change manager.\par
\textbf{Concerns}: What can change independently? Which interfaces or data migrations block progress? How is rollback achieved? \textbf{Aspects}: structural, behavioural and programmatic.\par
\textbf{Model kinds}: \texttt{MK-DEPENDENCY/1}, \texttt{MK-TRANSITION/1}. \textbf{Legends}: baseline/target, dependency strength, migration phase and cutover. \textbf{Correspondence method}: \texttt{CM-MIG/1}, each replaced component has a recorded target, retirement state and data/control handover.\par
\textbf{View methods}: change-impact traversal and migration gate review. \textbf{Sources}: lifecycle plan, inventory baseline and decision log. \textbf{Example view}: strangler migration sequence.\par
\subsection{VP-COST: Resource Economics viewpoint (version 1.0)}
\textbf{Purpose}: explain cost drivers and trade-offs among architecture options. \textbf{Perspectives}: financial, management and operations. \textbf{Stakeholders}: product owner, FinOps engineer, budget owner.\par
\textbf{Concerns}: Which elements dominate total cost? How do scaling assumptions affect cost? What is the consequence of a selected service topology? \textbf{Aspects}: economic and structural (local classification).\par
\textbf{Model kinds}: \texttt{MK-COST/1} (cost-allocation table and scenario model). \textbf{Legend}: units, currency, forecast horizon, estimated vs. measured value. \textbf{Correspondence method}: \texttt{CM-BILL/1}, each forecast cost item resolves to an identifiable architecture element and source assumption.\par
\textbf{View methods}: sensitivity analysis and alternative scenario comparison. \textbf{Sources}: vendor rate cards and operational usage measurements. \textbf{Example view}: service cost by workload.\par
\section{Model-kind specification templates \texorpdfstring{\Source{8.2}}{(8.2)}}
\begin{Requirement}\FlagReq{} For each model kind, include or reference modelling conventions; associated view and correspondence methods, where applicable; version identification under project/organization rules; and information sources. \Source{8.2}\end{Requirement}
\subsection{Full example: MK-DATAFLOW/1}
\begin{longtable}{p{0.23\linewidth}p{0.65\linewidth}}\toprule\textbf{Item}&\textbf{Specification for local use}\\\midrule\endhead
Identity&\texttt{MK-DATAFLOW}, version 1.0, domain software/data security.\\
Purpose&Express sources, processors, data stores and directed data exchanges for a defined scenario.\\
Entities&\texttt{ExternalParty}, \texttt{Process}, \texttt{DataStore}, \texttt{DataObject}, \texttt{TrustBoundary}.\\
Relationships&\texttt{DataFlow(source,target,dataObject)}, \texttt{Contains(boundary,entity)}.\\
Attributes&Unique ID, owner, sensitivity, environment, retention class, interface ID.\\
Constraints&A flow has two endpoints; every flow names an information payload; a boundary crossing carries a trust-relevant annotation. These are \FlagImpl{} local rules, not generic ISO requirements.\\
Representation&Arrows signify directed data movement; cylinders denote stores; outlined boxes denote processes. Legend must define any altered symbols.\\
Methods&\texttt{VM-DATA/1}: construct from interface inventory; \texttt{CM-DATA/1}: validate lineage coverage against classified data catalog.\\
Sources&Local modelling convention, reviewed data inventory, interface schema and approved privacy assessment.\\\bottomrule\end{longtable}
\subsection{Reusable model-kind form}
\Field{Identifier / version / scope / source}{[TODO]}
\Field{Syntax, semantic interpretation, constructs, relationships, constraints}{[TODO]}
\Field{Associated view methods and correspondence methods, if any}{[TODO]}
\Field{Examples, input sources and validation criteria}{[TODO]}
\section{Legend and view-method specifications}
\begin{Requirement}\FlagReq{} If a view component has no governing model kind, attach a legend defining its conventions. When required to explain construction or use, define view methods in the relevant specs. \Source{6.8, 8.1 and 8.3}\end{Requirement}
\begin{Worked}\textbf{Legend \texttt{LG-OPS-NARR/1}:} \texttt{MUST-ACT} identifies a mandatory local runbook step; \texttt{OWNER} records the response role; \texttt{SLO} is the service target; uncertainty is tagged \texttt{UNVERIFIED}. The legend tells readers how to interpret a \emph{narrative view component} rather than treating it as an ungoverned paragraph.\end{Worked}
\begin{longtable}{p{0.23\linewidth}p{0.65\linewidth}}\toprule\textbf{View method field}&\textbf{Example: VM-DATA/1}\\\midrule\endhead
Construction&Read approved interfaces; enumerate flows and stores; assign IDs and data classification; review with owners.\\
Interpretation&Arrow indicates data direction rather than call initiation; boundary shapes indicate trust context.\\
Analysis&Join flow records with classification and interface data; report unowned flows and missing risk treatment.\\
Input provenance&Inventory snapshot and schema revision; quality and freshness noted.\\
Output evidence&Revision-controlled view component and consistency-check result.\\\bottomrule\end{longtable}
\begin{Guidance}One model kind can be shared by several viewpoints and one view component can appear in several views. Always preserve unique identity and provenance; do not duplicate it without establishing correspondences or change-management rules.\end{Guidance}
\section{Catalog selection worksheet}
\begin{longtable}{p{0.15\linewidth}p{0.38\linewidth}p{0.30\linewidth}}\toprule\textbf{Viewpoint}&\textbf{Local concerns actually framed}&\textbf{Selected / evidence / exception}\\\midrule\endhead
VP-CXT&[TODO]&[TODO]\\ VP-FUN&[TODO]&[TODO]\\ VP-INFO&[TODO]&[TODO]\\ VP-DEP&[TODO]&[TODO]\\ VP-SEC&[TODO]&[TODO]\\ VP-OPS&[TODO]&[TODO]\\ VP-EVO&[TODO]&[TODO]\\ VP-COST&[TODO]&[TODO]\\\bottomrule\end{longtable}
\textbf{Selection rule (local):} choose viewpoints to cover the \emph{actual} concerns, perspectives and aspects identified in the AD, then make each selected specification and its governing components auditable. No viewpoint in this catalog is intrinsically required by the standard.
\LegalFooter
\end{document}
